\documentclass[a4paper,12pt]{article}
\usepackage{geometry}
\usepackage{graphicx}
\usepackage{setspace}

 \geometry{
 a4paper,
 total={170mm,257mm},
 left=20mm,
 top=20mm,
 }
\usepackage{verbatim}
\usepackage[utf8]{inputenc}
\usepackage{amsmath}
\usepackage{gensymb}
\usepackage{verbatim} % env for block comment 
\usepackage{ragged2e} % para usar flusleft
\usepackage[brazilian]{babel}
\usepackage{minted}
\usemintedstyle{tango}
\usepackage{amsfonts}


\begin{document}

\begin{center}
\thispagestyle{empty}
\onehalfspacing 
\begin{center} 
\includegraphics[scale=0.6]{logo.jpg} 
\end{center}

\vspace{2cm}
{\huge Álgebra Linear}\\
\vspace{2cm} % MUDAR ESTE VALOR PARA ALINHAR APÓS MUDAR A DESCRIÇÃO DA TAREFA

\noindent\rule{8cm}{0.4pt}
\vspace{10mm}

{\Large Tarefa 5}\\
\vspace{5mm}
% DESCRIÇÃO DA TAREFA
{ UNIDADE 4 - TRANSFORMAÇÕES LINEARES }\\
4.1 - Definição e propriedades.\\
4.2 - Transformações lineares no plano e no espaço.\\
4.3 - Núcleo e imagem de uma transformação linear.\\
4.4 - Matriz associada a uma transformação linear.\\
4.5 - Operações com transformações lineares.\\
4.6 - Operadores lineares e operadores inversíveis.\\
4.7 - Operadores ortogonais e simétricos.\\
\vspace{10mm}

\noindent\rule{8cm}{0.4pt}\\
\vspace{2cm} % MUDAR ESTE VALOR PARA ALINHAR APÓS MUDAR A DESCRIÇÃO DA TAREFA

{\large Prof. José Antônio Trindade Borges da Costa}\\
Tutora: Dionéia Migotto\\
Aluno: Tiago Cardoso Floriano\\
Matrícula: 202021740 %\hline
\end{center}

\newpage
\onehalfspacing 
\noindent
\textbf{ATIVIDADE 1}, 5) Dada a função $T:\mathbb{R}^2 \rightarrow \mathbb{R}^3$, definida por T(x, y) = (x+ky, x+k, y), verificar para quais dos seguintes valores de $k$, $T$ é uma transformação linear: \\
\\
a. $k = x$ \\
\\
{\textbf{SOLUÇÃO:}} \\
Tomando o exemplo 1 (p. 142) \cite{ines}, e considerando $\vec{u}=(x_1,y_1)$ e $\vec{v}=(x_2,y_2)$, teremos em $T(\vec{u}+\vec{v})$:\\
\begin{equation}
    \begin{split}
        T(\vec{u}+\vec{v})=T(x_1+x_2, y_1+y_2)=\\
        ((x_1+x_2)+(x_1+x_2)(y_1+y_2),(x_1+x_2)+(x_1+x_2), (y_1+y_2))=\\
        ((x_1+x_2)+(x_1 y_1)+(x_2 y_2),(2x_1+2x_2),(y_1+y_2))=\\
        ((2x_1 y_1+2x_2 y_2), (2x_1 +2x_2), (y_1+y_2))=\\
        ((2x_1 y_1, 2x_1, y_1) + (2x_2 y_2, 2x_2, y_2)=\\
        T(x_1,y_1)+T(x_2,y_2)=\\
        T(\vec{u}+\vec{v})
    \end{split}
\end{equation}



\newpage
\noindent 
\bibliographystyle{abbrv}
\begin{thebibliography}{9}
%... items da bibliografia aqui

\bibitem{ines} FERREIRA, Inês Farias,
	\emph{Álgebra Linear}, UAB-UFSM, 2010.
\bibitem{leon} LEON, Steven J.,
	\emph{Álgebra Linear com Aplicações}, 9th ed. Rio de Janeiro: LTC; 2018. ISBN: 978-85-216-3577-2.

\end{thebibliography}
 
\end{document}